\begin{afterword}
\summaryhead

The post defines the affect heuristic as the use of quick feelings of good or bad as a basis for judgment, and describes a series of studies. People would pay more than twice as much to insure a clock they love, though the policy pays the same \$100 for either clock. They judged a disease that kills 1,286 people out of 10,000 more dangerous than one with a 24.14\% chance of death. They supported an airport safety measure more when it would save 98\% of 150 lives than when it would save 150 lives. And they often preferred to draw from a bowl with 7 winning beans in 100 rather than 1 in 10, while knowing the odds were worse.

Finucane and colleagues found that information about a technology's benefits lowered its perceived risk, and the reverse, with a stronger effect under time pressure. Ganzach found that analysts judged unfamiliar stocks as all good or all bad, while for familiar stocks they judged risk and return to rise together. People merge the separate good and bad aspects of a thing into one feeling.

\responsehead

The post is a survey of studies from the work of Slovic and colleagues, and most of it is reported accurately. The airport figures match the published table. The jelly-bean study matches its abstract, including the subjects' report that they knew the odds were against them. And this time the Finucane study is described by its real design, where ``The Scales of Justice, the Notebook of Rationality'' had presented an example of its own as the finding. The problems are in a few details and in one explanation.

The clock study is presented as a bias that one might excuse as ``purchase of consolation.'' That is the explanation its authors gave. Their paper proposes a ``consolation hypothesis,'' in which people treat insurance compensation as a token of consolation. The post presents the authors' account as a possible excuse and does not say it is theirs.

One of Finucane's results is stated more strongly than the original paper states it. The post says time pressure ``greatly'' increased the effect, as the authors' later summaries also say, but in the 2000 paper the effect is modest. The mean correlation between risk and benefit ratings was $-0.45$ under time pressure against $-0.33$ without it, a difference that approached significance.

The disease study is reworded. In the summaries I could read, both descriptions were death rates in a population, and later work classes the result with the ratio bias: the larger count of deaths looks worse. The post describes the second figure as a chance of death and explains the result by a thousand bodies against ``a single person who's more likely to survive than not.'' The explanation is hedged with ``Apparently,'' but it depends on the rewording. I could read only summaries of this study and of the clock study.

Finally, the article recommended for further reading covers only two of the post's studies. The others are summarized in the same authors' ``The Affect Heuristic.''

In short: an accurate survey of the affect-heuristic research for most of its length, which offers the clock study's own explanation without saying it is the authors', repeats the authors' later overstatement of one of Finucane's results, and rewords the disease study in a way its explanation depends on.
\end{afterword}
